% BHU PYQ -- Manually transcribed & error-corrected
% Paper: BCH-122 : Financial Accounting-II
% Exam: B.Com. (Hons.) Semester II Examination, 2021-22

\documentclass[11pt,a4paper]{article}
\usepackage[a4paper,top=2.5cm,bottom=2.5cm,left=2.8cm,right=2.8cm,headheight=14pt]{geometry}
\usepackage{amsmath,amssymb}
\usepackage[T1]{fontenc}
\usepackage[utf8]{inputenc}
\usepackage{lmodern,microtype,fancyhdr,enumitem,hyperref,lastpage,parskip}
\usepackage{booktabs}

\hypersetup{
    colorlinks=true,
    urlcolor=black,
    linkcolor=black,
    pdftitle={BCH-122: Financial Accounting-II -- BHU B.Com. (Hons.) Sem II 2021-22}
}

\pagestyle{fancy}
\fancyhf{}
\renewcommand{\headrulewidth}{0.4pt}
\renewcommand{\footrulewidth}{0.4pt}
\fancyhead[L]{\small   \textit{Banaras Hindu University}}
\fancyhead[R]{\small   \textit{BCH-122: Financial Accounting-II}}
\fancyfoot[C]{\small Page  \thepage\ of \pageref{LastPage}}


\newlist{parts}{enumerate}{2}
\setlist[parts,1]{label=(\alph*),leftmargin=*,itemsep=5pt,topsep=3pt}
\setlist[parts,2]{label=(\roman*),leftmargin=*,itemsep=3pt,topsep=2pt}

\newcommand{\pts}[1]{\hfill   \textit{[#1]}}


\begin{document}


\begin{center}
  {\large\bfseries BANARAS HINDU UNIVERSITY}\\[2pt]
  {
\normalsize B.Com. (Hons.) --- Semester II Examination, 2021-22}\\[6pt]
  \rule{0.8\linewidth}{0.5pt}\\[6pt]
  {\large\bfseries Paper No. BCH-122: Financial Accounting-II}\\[4pt]
  {
\normalsize Time: 3 Hours\hfill Maximum Marks: 70}
\end{center}

\rule{\linewidth}{0.4pt}

\smallskip



\noindent   \textit{Instructions: Answer all questions. All questions carry equal marks (14 marks each).}

\bigskip




\noindent   \textbf{Question 1.}
Differentiate between Promissory Note and Bills of Exchange. Give Journal Entries in the Books of Drawer and Drawee.\pts{14}

\centerline{   \textbf{OR}}


\begin{parts}
  \item State the causes of difference between the balance of Cash Book and Pass Book.\pts{6}
  \item From the following particulars, prepare Bank Reconciliation Statement to ascertain the balance as per Cash Book:\pts{8}
  
\begin{itemize}[itemsep=2pt,topsep=2pt]
    \item[(i)] Credit balance of Pass Book: Rs. 12,000
    \item[(ii)] Cheque issued but not presented: Rs. 3,000
    \item[(iii)] Bank debited Rs. 50 and credited Rs. 170
    \item[(iv)] Received a cheque and deposited: Rs. 350
    \item[(v)] Cheques deposited but not credited: Rs. 1,000
    \item[(vi)] Insurance premium paid by bank: Rs. 500
    \item[(vii)] Bank credited interest on investment: Rs. 250
  \end{itemize}
\end{parts}

\medskip\rule{\linewidth}{0.2pt}


\newpage




\noindent   \textbf{Question 2.}
AB Ltd. issued a prospectus equity application for 2000 shares of Rs. 10 each at a premium of Rs. 2 per share, payable as Rs. 2 on application, Rs. 5 on allotment, Rs. 3 on first call and Rs. 2 on second call. Applications were received for 3400 shares and prorata allotment was made on application for 2400 shares and remaining applications were rejected. Money overpaid on applications were adjusted on allotment. Mr. Rakesh to whom 40 shares were allotted failed to pay allotment and first call and his shares were forfeited. Mr. Manish the allottee of 60 shares failed to pay the two calls and his shares were forfeited. All these forfeited shares were reissued as fully paid for Rs. 8 per share. Journalise the above transactions.\pts{14}

\centerline{   \textbf{OR}}

Discuss the conditions for redemption of preference shares. Give the various journal entries to be passed on redemption of preference shares.\pts{14}

\medskip\rule{\linewidth}{0.2pt}




\noindent   \textbf{Question 3.}
What do you mean by Sinking Fund Method? Give the various journal entries and show necessary accounts to redeem debentures under Sinking Fund Method.\pts{14}

\centerline{   \textbf{OR}}


\begin{parts}
  \item XY Ltd. purchased machinery worth Rs. 5,00,000 from Z Ltd. and payment was made by issue of 10\% debentures of Rs. 100 each. Pass necessary Journal entries for the purchase of machinery and issue of debentures when:
  
\begin{enumerate}[label=(\roman*),itemsep=2pt]
    \item debentures are issued at par
    \item debentures are issued at 10\% discount
    \item debentures are issued at 10\% premium
  \end{enumerate}\pts{7}
  \item Journalise the following transactions:\pts{7}
  
\begin{enumerate}[label=(\roman*),itemsep=2pt]
    \item Debentures issued at Rs. 100 to be redeemed at Rs. 120
    \item Debentures issued at Rs. 90 to be redeemed at Rs. 105
    \item Debentures issued at Rs. 95 to be redeemed at Rs. 100
    \item 700 debentures were issued at face value of Rs. 100
  \end{enumerate}
\end{parts}

\medskip\rule{\linewidth}{0.2pt}


\newpage




\noindent   \textbf{Question 4.}
Q Ltd. was registered to take over a running business from 1st April, 2021. The company was incorporated on 1st August, 2021 and certificate to commence business was received on 1st October, 2021. The company prepared the following Profit \& Loss Account on 31st March, 2022:

\begin{center}
  \small
  
\begin{tabular}{lr|lr}
     oprule
   \textbf{Particulars} &   \textbf{Amount (Rs.)} &   \textbf{Particulars} &   \textbf{Amount (Rs.)} \\
    \midrule
    To Rent & 20,000 & By Gross Profit & 2,00,000 \\
    To Insurance & 10,000 & & \\
    To Electricity & 7,400 & & \\
    To Salaries & 51,000 & & \\
    To Director fees & 5,600 & & \\
    To Commission & 7,000 & & \\
    To Advertisement & 4,000 & & \\
    To Office expenses & 7,500 & & \\
    To Depreciation & 3,000 & & \\
    To Preliminary exp. & 7,500 & & \\
    To Bad debts & 2,000 & & \\
    To Interest & 3,000 & & \\
    To Net Profit & 72,000 & & \\
    \midrule
   \textbf{Total} &   \textbf{2,00,000} &   \textbf{Total} &   \textbf{2,00,000} \\
    ottomrule
  \end{tabular}
\end{center}




\noindent   \textbf{Additional Information:}

\begin{parts}
  \item Sales up to commencement of business amounted to Rs. 3,00,000 wherefore it spurted to record an increase of two-fifth during the rest of the year.
  \item Depreciation included Rs. 1,200 for the assets acquired in post-incorporation period.
\end{parts}
Prepare Profit \& Loss Account to show profit prior to and post-incorporation indicating the basis of allocation.\pts{14}

\centerline{   \textbf{OR}}

T Ltd, a public company was incorporated on 1st March, 2020 and received its certificate to commence business on 1st April, 2020. It took over the business of S Ltd. running from 1.1.2020. From the following particulars, prepare Profit \& Loss Account to show profit prior to and post-incorporation:

\begin{center}
  \small
  
\begin{tabular}{lr|lr}
     oprule
   \textbf{Particulars} &   \textbf{Amount (Rs.)} &   \textbf{Particulars} &   \textbf{Amount (Rs.)} \\
    \midrule
    Salaries & 20,000 & Preliminary expenses & 6,000 \\
    Rent & 28,000 & Bad debts & 900 \\
    Director fees & 3,500 & Discount & 3,600 \\
    Interest on debentures & 8,000 & Interest to vendor & 10,000 \\
    Auditor fees & 12,500 & Printing and stationery & 5,200 \\
    Depreciation & 12,000 & Commission & 4,800 \\
    Advertising & 24,000 & Establishment charges & 11,400 \\
    ottomrule
  \end{tabular}
\end{center}
Total sales for the year was Rs. 9,00,000 while sales up to 31.03.2020 was Rs. 3,00,000. The Gross Profit for the year was Rs. 2,49,000.\pts{14}

\medskip\rule{\linewidth}{0.2pt}


\newpage




\noindent   \textbf{Question 5.}
The following is the Trial Balance of S. S. Ltd. as at 31st March, 2021. Prepare Trading and Profit \& Loss Account for the year and Balance Sheet as on 31st March, 2021:

\begin{center}
  \small
  
\begin{tabular}{lrr}
     oprule
   \textbf{Particulars} &   \textbf{Debit Balance (Rs.)} &   \textbf{Credit Balance (Rs.)} \\
    \midrule
    Land and Building & 34,000 & \\
    Plant and Machinery & 66,000 & \\
    Fees & 2,000 & \\
    Calls in Arrear & 600 & \\
    5\% Government Bonds & 3,600 & \\
    Motor Vehicle & 4,000 & \\
    Interim Dividend & 1,800 & \\
    Goodwill & 54,600 & \\
    Repairs & 300 & \\
    Purchases & 96,000 & \\
    Interest to Bank & 840 & \\
    Audit Fees & 400 & \\
    Advertisement & 1,000 & \\
    Carriage & 1,500 & \\
    Wages & 9,200 & \\
    Insurance & 2,000 & \\
    Opening Stock & 19,000 & \\
    Tools Purchased on Instalment & 4,000 & \\
    Loss on Foreign Exchange & 1,700 & \\
    Furniture & 2,900 & \\
    Trade Debtors & 8,300 & \\
    Return Inward & 2,800 & \\
    Share Capital & & 80,000 \\
    Global Depository Receipts & & 12,000 \\
    Pension Fund & & 6,000 \\
    P \& L A/c Balance & & 3,540 \\
    Return Outwards & & 2,000 \\
    Sales & & 1,23,000 \\
    Suppliers of Material & & 10,000 \\
    Private Placement of Debentures & & 30,000 \\
    Income Tax Outstanding & & 5,000 \\
    Foreign Exchange Credit & & 40,000 \\
    Creditors & & 5,000 \\
    \midrule
   \textbf{Total} &   \textbf{3,16,540} &   \textbf{3,16,540} \\
    ottomrule
  \end{tabular}
\end{center}




\noindent   \textbf{Adjustments:}

\begin{parts}
  \item Closing stock Rs. 22,000.
  \item Amount spent Rs. 5,600 on protection of environment not yet accounted for.
  \item Audit fees include Rs. 100 for tax consultancy.
  \item Managing Director's commission to be provided at 2\% on net profit before charging such commission.
  \item An amount of Rs. 2,000 to be written-off for Goodwill.
  \item Running expenses for the previous year of Rs. 400 not provided in that year, paid now in current year.
\end{parts}
\pts{14}

\centerline{   \textbf{OR}}

Explain the provisions of Companies Act as regards the preparation of Profit and Loss Account.\pts{14}

\vfill
\rule{\linewidth}{0.4pt}

\begin{center}\small
\href{https://bhu-exam.vercel.app/}{Click here to visit our website}\\[4pt]
\href{https://me-aryan.vercel.app/}{Click here to visit our contributor}\\[4pt]
\href{https://quashan-me.vercel.app/}{Click here to visit our contributor}
\end{center}

\end{document}
